\chapter{Селектор source-линии в семейно замкнутой среде}
\label{ch:v8-baryon-c0-so3-closed-environment-source-line-selector-gate}

Предыдущий гейт оставил семейство изометрических интертвейнеров,
параметризованное $[u:v]\in\mathbb {RP}^1$. Однако физический коннектор
должен одновременно быть нечётным относительно уже фиксированных
градуировок. На трёх target- и двух source-направлениях они равны
\begin{equation}
 \Gamma_{\rm tar}=\operatorname{diag}(-1,+1,+1),
 \qquad \Gamma_{\rm src}=\operatorname{diag}(+1,-1).
 \label{eq:v8-baryon-c0-source-line-current-gradings}
\end{equation}
Именно это разложение допускает текущие три стрелки:
\begin{equation}
 \Gamma_{\rm tar}E_{00}+E_{00}\Gamma_{\rm src}=0,
 \quad
 \Gamma_{\rm tar}E_{11}+E_{11}\Gamma_{\rm src}=0,
 \quad
 \Gamma_{\rm tar}E_{21}+E_{21}\Gamma_{\rm src}=0.
 \label{eq:v8-baryon-c0-current-three-arrows-odd}
\end{equation}
Но все три стрелки, необходимые для $SO(3)$-замыкания, имеют ненулевой
остаток нечётности:
\begin{equation}
 \left\|\Gamma_{\rm tar}E_{ij}+E_{ij}\Gamma_{\rm src}\right\|_{\rm HS}^2=4,
 \qquad E_{ij}\in\{E_{10},E_{20},E_{01}\}.
 \label{eq:v8-baryon-c0-missing-arrows-even-residual}
\end{equation}
Причина структурна: семейное действие не коммутирует с target-grading,
\begin{equation}
 \bigl(\|[\Gamma_{\rm tar},J_1]\|_{\rm HS}^2,
 \|[\Gamma_{\rm tar},J_2]\|_{\rm HS}^2,
 \|[\Gamma_{\rm tar},J_3]\|_{\rm HS}^2\bigr)=(0,8,8).
 \label{eq:v8-baryon-c0-target-grading-so3-commutator}
\end{equation}
По лемме Шура неприводимая тройка могла бы иметь только скалярную
градуировку, тогда как текущая имеет спектральный тип $1+2$.

\section{Ни одна точка RP1 не является нечётной}

Для $z\in\mathbb R^3$ общий семейный интертвейнер записывается как
\begin{equation}
 T_{u,v}(z)=\begin{pmatrix}uz&vz\end{pmatrix},
 \qquad u^2+v^2=1.
 \label{eq:v8-baryon-c0-source-line-general-family-map}
\end{equation}
Условие нечётности для всех трёх базисных $z$ даёт линейную систему
\begin{equation}
 \Gamma_{\rm tar}T_{u,v}(z)+T_{u,v}(z)\Gamma_{\rm src}=0,
 \qquad \rank\mathcal C_{\rm odd}=2,\qquad (u,v)=(0,0).
 \label{eq:v8-baryon-c0-source-line-oddness-kernel-zero}
\end{equation}
Следовательно, множество допустимых проективных точек пусто. Суммарный
точный дефект равен
\begin{equation}
 \Delta_{\Gamma}(u,v)=
 \sum_{k=1}^3\|\Gamma_{\rm tar}T_{u,v}(e_k)
 +T_{u,v}(e_k)\Gamma_{\rm src}\|_{\rm HS}^2
 =8u^2+4v^2.
 \label{eq:v8-baryon-c0-source-line-parity-defect}
\end{equation}
На двух координатных линиях
\begin{equation}
 \Delta_{\Gamma}(1,0)=8,\qquad
 \Delta_{\Gamma}(0,1)=4,\qquad
 \min_{u^2+v^2=1}\Delta_{\Gamma}=4>0.
 \label{eq:v8-baryon-c0-source-line-axis-defects}
\end{equation}
Даже проекция на нечётную часть не восстанавливает тройку:
\begin{equation}
 \rank P_{\rm odd}T_{1,0}=1,\qquad
 \rank P_{\rm odd}T_{0,1}=2.
 \label{eq:v8-baryon-c0-source-line-odd-projection-ranks}
\end{equation}
То есть меньший дефект линии $a_0$ не является точным допуском.

\section{Внутренние кандидаты селектора}

Нормированный след и паулиевский пакет нового $M_2$-блока изотропны:
\begin{equation}
 \rho_\tau=\frac12I_2,\qquad
 X^2+Y^2+H^2=3I_2,\qquad
 \Tr(\rho_\tau P_s)=\Tr(\rho_\tau P_a)=\frac12.
 \label{eq:v8-baryon-c0-source-line-trace-pauli-isotropy}
\end{equation}
Real-структура сохраняет обе координатные линии, а деполяризующая динамика
имеет ту же стационарную матрицу $I_2/2$. Единственный доступный
анизотропный оператор, сама source-grading, даёт семейство
\begin{equation}
 h_{\alpha,\beta}=\beta I_2+\alpha\Gamma_{\rm src},
 \qquad \Spec h_{\alpha,\beta}=\{\beta+\alpha,\beta-\alpha\},
 \label{eq:v8-baryon-c0-source-line-grading-hamiltonian}
\end{equation}
но знак $\alpha$ не выведен, а отождествление grading с энергией является
новым parent-постулатом.

Формально можно заменить target-grading скаляром:
\begin{equation}
 \begin{aligned}
 \Gamma_{\rm tar}^{(-)}=-I_3&\Rightarrow [u:v]=[1:0]
 &&\text{и две смены знака},\\
 \Gamma_{\rm tar}^{(+)}=+I_3&\Rightarrow [u:v]=[0:1]
 &&\text{и одна смена знака}.
 \end{aligned}
 \label{eq:v8-baryon-c0-source-line-uniform-grading-extensions}
\end{equation}
Минимум числа переприсвоенных знаков предпочитает второй вариант, но это
критерий редактирования архитектуры, а не следствие действия; кроме того,
он меняет grading уже существующего состояния $H_{21}$.

Итоговый реестр селекторов равен
\begin{equation}
 N_{\rm source\ selector}=0/6:
 \quad\{\Gamma_{\rm src},\tau_2,\text{Pauli-Casimir},
 \mathcal L_{45},J_0,h_{\alpha,\beta}\}.
 \label{eq:v8-baryon-c0-source-line-selector-ledger}
\end{equation}

\begin{theorem}[Градуировочный запрет сильнее неоднозначности $\mathbb {RP}^1$]
Ни одна точка семейства $[u:v]\in\mathbb {RP}^1$ не задаёт ненулевой
$SO(3)$-эквивариантный нечётный коннектор при текущих градуировках.
Следовательно, задача не сводится к выбору линии: сначала требуется новый
семейный target-триплет с однородной градуировкой.
\end{theorem}

\begin{nogocriterion}[Минимальный parity-дефект не является допуском]
Линия $[0:1]$ минимизирует $\Delta_\Gamma$, но минимум равен $4$, а не
нулю. Нельзя подменять точное условие нечётности наименьшим нарушением и
объявлять эту линию выбранной картой $c_0$.
\end{nogocriterion}

\begin{protocol}\begin{enumerate}
\item текущая target-grading: \textbf{$(-,+,+)$};
\item текущая source-grading: \textbf{$(+,-)$};
\item недостающие стрелки нечётны: \textbf{нет};
\item нечётный family-Hom: \textbf{$0$};
\item допустимых точек $\mathbb {RP}^1$: \textbf{$0$};
\item минимальный parity-дефект: \textbf{$4$};
\item ранги нечётных проекций осей: \textbf{$1,2$};
\item внутренних селекторов: \textbf{$0/6$};
\item необходимое расширение: \textbf{однородно градуированный triplet};
\item следующий гейт: \textbf{endpoint-допуск grading-совместимого family-triplet}.
\end{enumerate}\end{protocol}

Машинный сертификат сохранён в
\path{s2t_v8_baryon_c0_so3_closed_environment_source_line_selector_gate_results.json}.